\section{Background}
\label{sec:background}

\begin{figure*}[!t]
  \centering
  \definecolor{csScaf}{RGB}{225,225,225}
  \definecolor{csTask}{RGB}{204,224,245}
  \definecolor{csDesc}{RGB}{207,236,214}
  \definecolor{csVal}{RGB}{252,224,195}
  \definecolor{csSchema}{RGB}{230,214,240}
  \setlength{\fboxsep}{1.5pt}
  \newcommand{\ph}[1]{{\normalfont\itshape$\langle$#1$\rangle$}}
  \newcommand{\csblock}[2]{\par\noindent\colorbox{#1}{\parbox{\dimexpr\linewidth-2\fboxsep\relax}{\raggedright\ttfamily\tiny #2}}\par\vspace{1pt}}
  \newcommand{\csrole}[1]{\par\noindent{\sffamily\tiny\color{black!55}#1}\par\vspace{0.3pt}}
  \newcommand{\cshead}[1]{\par\noindent{\sffamily\scriptsize\bfseries #1}\par\vspace{1.5pt}}
  \newcommand{\cskey}[2]{\fcolorbox{black!40}{#1}{\rule{0pt}{1.1ex}\rule{1.1ex}{0pt}}\,{\sffamily\tiny #2}}
  \cskey{csScaf}{framework scaffolding}\hspace{1em}%
  \cskey{csTask}{task instruction}\hspace{1em}%
  \cskey{csDesc}{parameter descriptions}\hspace{1em}%
  \cskey{csVal}{argument values}\hspace{1em}%
  \cskey{csSchema}{output schema}
  \vspace{3pt}

  \begin{minipage}[t]{0.31\textwidth}
    \cshead{byLLM}
    \csrole{system}
    \csblock{csScaf}{This is a task you must complete by returning only the output. ...}
    \csrole{user}
    \csblock{csTask}{assess\_evidence(search\_results: str, question: str, claim: str, evidences: list[Evidence]) -> Finding --- \ph{sem}}
    \csblock{csDesc}{search\_results: str ---- \ph{sem}\newline question: str ---- \ph{sem}\newline claim: str ---- \ph{sem}\newline evidences: list[Evidence] ---- \ph{sem}}
    \csblock{csVal}{search\_results = \ph{value}\newline question = \ph{value}\newline claim = \ph{value}\newline evidences = \ph{value}}
    \csblock{csSchema}{Schema requirements: \ph{Finding fields}}
    \csrole{response\_format}
    \csblock{csSchema}{\ph{Finding JSON schema}}
  \end{minipage}\hfill
  \begin{minipage}[t]{0.31\textwidth}
    \cshead{DSPy}
    \csrole{system}
    \csblock{csDesc}{Your input fields are:\newline 1. `search\_results` (str): \ph{desc}\newline 2. `question` (str): \ph{desc}\newline 3. `claim` (str): \ph{desc}\newline 4. `evidences` (list[Evidence]): \ph{desc}\newline Your output fields are: 1. `finding` (Finding): \ph{desc}}
    \csblock{csScaf}{All interactions will be structured ...\newline [[ \#\# search\_results \#\# ]] \{search\_results\}\newline ... [[ \#\# evidences \#\# ]] \{evidences\}}
    \csblock{csSchema}{[[ \#\# finding \#\# ]] \{finding\} \# note: ... JSON schema: \ph{Finding JSON schema}}
    \csblock{csScaf}{[[ \#\# completed \#\# ]]}
    \csblock{csTask}{In adhering to this structure, your objective is: \ph{docstring}}
    \csrole{user}
    \csblock{csVal}{[[ \#\# search\_results \#\# ]] \ph{value}\newline [[ \#\# question \#\# ]] \ph{value}\newline [[ \#\# claim \#\# ]] \ph{value}\newline [[ \#\# evidences \#\# ]] \ph{value}}
    \csblock{csScaf}{Respond with the corresponding output fields, starting with `[[ \#\# finding \#\# ]]` ...}
  \end{minipage}\hfill
  \begin{minipage}[t]{0.31\textwidth}
    \cshead{NVIDIA nooa}
    \csrole{system}
    \csblock{csScaf}{<strategy\_prompt> You are generating structured data matching the specified return type. ... </strategy\_prompt>}
    \csrole{user}
    \csblock{csScaf}{<sys tag="1"> Task(prompt='''\# Your task}
    \csblock{csTask}{\ph{docstring}}
    \csblock{csDesc}{Args:\newline \hspace*{1em}search\_results: \ph{arg doc}\newline \hspace*{1em}question: \ph{arg doc}\newline \hspace*{1em}claim: \ph{arg doc}\newline \hspace*{1em}evidences: \ph{arg doc}}
    \csblock{csVal}{\#\# Input parameters:\newline search\_results = \ph{value}\newline question = \ph{value}\newline claim = \ph{value}\newline evidences = \ph{value}}
    \csblock{csScaf}{Perform the task and return the result directly.''') </sys>}
    \csrole{response\_format}
    \csblock{csSchema}{\ph{Finding JSON schema}}
  \end{minipage}
  \par\smallskip
  {\raggedright\sffamily\tiny
  \ph{sem}, \ph{docstring}, \ph{desc} and \ph{arg doc} are the description strings each framework takes from the program (byLLM \texttt{sem} strings, DSPy Signature docstring and field \texttt{desc}, nooa function docstring and its \texttt{Args} entries); \ph{value} is the runtime argument.\par}
  \caption{The prompt structure of LLM invocation function \texttt{assess\_evidence}, lowered by three frameworks. They have similar structures: Same five kinds of component, only relative position differs.}
  \label{fig:common-structure}
\end{figure*}


\subsection{Serving systems for LLM agents}
\label{sec:background-serving}

LLM serving engines retain the attention keys and values (KV) of tokens they have already processed to avoid repeating computation. PagedAttention manages these states in shareable memory blocks~\cite{kwon2023pagedattention}, while SGLang indexes cached token sequences in a radix tree~\cite{zheng2024sglang}. An incoming request reuses the longest cached prefix of its prompt and computes KV for the remaining suffix. Reuse stops at the first token mismatch: that token and every subsequent token must be recomputed, even if identical text appears later in a cached sequence, because its preceding context has changed. Prompt Cache supports modular reuse beyond a shared prefix, but requires modules to retain fixed, schema-assigned token positions~\cite{gim2024promptcache}.

Retained KV competes with running requests for device memory. Under memory pressure, engines such as SGLang evict unused cached prefixes by recency~\cite{zheng2024sglang}. Hierarchical caching systems, including CachedAttention~\cite{gao2024cachedattention}, Mooncake~\cite{qin2025mooncake}, and SGLang HiCache~\cite{sglang2026hicache}, extend the cache beyond device memory with host memory and, optionally, storage. KV evicted from the device can remain in the host tier and be transferred back when needed; prefetching can overlap this transfer with other work. Reloading retained KV can be cheaper than recomputing it. The host tier is also limited: admitting one sequence consumes space that could retain another, so offloading does not eliminate cache-capacity pressure.

An agent \emph{session} is one execution of an agent task, comprising a series of dependent LLM calls separated by tool and code execution. Between calls, the session's KV sits idle and is exposed to recency eviction precisely while the agent is doing work that may lead to its reuse. Agent-aware serving systems use execution information to manage these gaps. KVFlow uses an agent step graph to estimate steps to future invocations and guide eviction and prefetching~\cite{pan2025kvflow}. Continuum retains KV across tool calls with a time-to-live~\cite{li2025continuum}. CacheScout learns transitions between agents online to guide retention and prefetching~\cite{zhang2026cachescout}, while PBKV predicts future agent invocations from workflow history and the current context~\cite{zheng2026pbkv}. These systems share a request- or agent-step-level view of future execution: they organize cache management around the reuse of whole requests or agent steps.

\subsection{Agent frameworks with LLM calls as programming abstractions}
\label{sec:background-frameworks}

A class of agent frameworks exposes model-backed computation through typed programming interfaces. Examples include byLLM's meaning-typed programming (MTP) abstraction~\cite{dantanarayana2025mtp}, DSPy's Signatures and Modules~\cite{khattab2024dspy}, and NOOA's typed Python methods~\cite{furgale2026nooa}. SIGIL likewise compiles skill descriptions into harnesses containing typed model operations~\cite{dantanarayana2026sigil}. The developer declares a function or signature describing the desired behavior, its inputs, and its return type; invoking it is an ordinary call within the program's control flow. For example, \texttt{assess\_evidence} in Figure~\ref{fig:heterogeneity-sample} accepts search results, a question, a claim, and accumulated evidence, and returns a typed finding. We use \emph{call site} to mean a location in the program that invokes such an LLM-backed operation. A session is a concrete execution through these sites; a loop can visit the same site repeatedly with different argument values.

At each invocation, the framework runtime lowers the declaration and its argument values into model-facing messages. The declaration supplies the function name, docstring or semantic strings, parameter types and descriptions, and return schema. This information is known when the program is compiled: byLLM captures it in MTP's meaning-based intermediate representation, MT-IR~\cite{dantanarayana2025mtp}. The argument values become available only at runtime, when the call executes. Thus, a generated prompt combines a description fixed by the program with values supplied by that execution. 

Some frameworks make deliberate choices about how to arrange this information. NOOA places static context first, appended event history next, and volatile state last to maximize KV-cache reuse across turns~\cite{furgale2026nooa}. DSPy's adapters use a fixed ordering of system instructions, demonstrations, conversation history, and current inputs; its ChatAdapter emits input fields in signature order~\cite{dspy2026chatadapter}. These layouts follow hand-chosen conventions applied across call sites and programs, rather than an order derived from the properties of a particular program's bindings. The framework knows the structure behind every prompt, but the serving engine receives only the token sequence.
